%%%%%%%%%%%%%%%%%%%%%%%%%%%%%%%%%%%%%%%%%%%%%%%%%%%%%
%
% This is the template for the EGAS 55 conference abstracts.
% Please, insert your name and affiliations, text of the abstract and references
% at the locations in the file where they are mentioned.
% Do not change the part defining the format of the file.
%%%%%%%%%%%%%%%%%%%%%%%%%%%%%%%%%%%%%%%%%%%%%%%%%%%%%
%
%%%%%%%%%%%%%%%%%%%%   Format, do not change  %%%%%%%%%%%%%%%%%
%

% LaTeX file to be processed by pdflatex

\documentclass[11pt,a4paper]{article}

\usepackage{times}
\usepackage{graphicx}

%\usepackage{draftwatermark}


\setlength{\hoffset}{0.0cm}
\setlength{\voffset}{0.0cm}

\setlength{\oddsidemargin}{-0.8cm}
\setlength{\marginparwidth}{-0.5cm}

\setlength{\headsep}{0.0cm}
\setlength{\headheight}{0.0cm}
\setlength{\topskip}{0.4cm}
\setlength{\topmargin}{-0.5cm}
\setlength{\footskip}{1.0cm}

\setlength{\textwidth}{16.9cm}
\setlength{\textheight}{24.7cm}

\setlength{\parindent}{0em}


\newcommand{\abstracttitle}[1]{\parbox{\textwidth}
                              {\setlength{\baselineskip}{19pt}%
                               \begin{center}{\large\bfseries #1}\end{center}}}
\newcommand{\authors}[1]{\begin{center}#1\end{center}}
\newcommand{\addresses}[1]{\begin{center}#1\end{center}}

\renewcommand{\labelenumi}{[\arabic{enumi}]}

\pagestyle{empty}
% \pagestyle{headings}
\markright{EGAS 55}

\usepackage{siunitx}
\usepackage{biblatex}
\bibliography{USOC.bib}
% \bibliographystyle{plainnat}

%%%%%%%%%%%%%%   End format   %%%%%%%%%%%%%%%%%


%%%%%%%%%%%%%%%%%%%%%%%%%%%%%%%%%%%%%%%%%%%
%             Editable part
%%%%%%%%%%%%%%%%%%%%%%%%%%%%%%%%%%%%%%%%%%%
%

\begin{document}
%
% Title of the abstract
\abstracttitle{USOC: Ultra-precise, shock-resistant optical clocks}
%
% Authors and affiliation
\authors{\textit{\underline{Thomas Walker*}, Ludovico Iannizzotto Venezze, Leonie Hawkins, Charles Baynham, Richard Hobson, Oliver Buchmueller}}

\addresses{{\small Imperial College London, United Kingdom\\ \vspace{0.2cm}
			(*) t.walker@imperial.ac.uk}}

\setlength{\parskip}{\bigskipamount}
%%%%%%%%%%%%%%%%%%%%%%%%%%%%%%%
% Text of your abstract
%%%%%%%%%%%%%%%%%%%%%%%%%%%%%%%

Optical atomic clocks require atoms to be cooled down to ultracold temperatures. 
In optical lattice clocks, this imposes long dead times between clock interrogations, 
leading to the Dick effect - aliasing of high-frequency local oscillator noise into low-frequency clock noise. 
This can be circumvented by using a continuous scheme, 
in which atoms are cooled and loaded into a moving optical lattice which conveys them between spatially separated interrogation regions \cite{katori_longitudinal_2021}. 
Dick noise is the the most significant source of noise in current neutral atom optical clocks, 
and eliminating it would potentially allow these clocks to reach the quantum projection noise limit \cite{schioppo_ultrastable_2017}. 

This talk will introduce a new strontium-based continuous optical lattice clock, 
loaded from a MOT operating on the mid-infrared \SI{2.92}{\micro\meter} $5s5p^3P_2$ to $5s4d^3D_3$ transition. 
Operating the MOT on these metastable states has several advantages over the typical 689 nm MOT used for Sr; most significantly, 
it can operate close to the interrogation region without significantly perturbing the 698nm 5s2 1S0 to 5s5p 3P0 clock transition. 
Additionally, the 689 nm MOT  requires additional lasers for population recycling,
 and has a higher recoil limit (\SI{250}{\nano\kelvin} vs \SI{12}{\nano\kelvin}) \cite{hobson_midinfrared_2020, akatsuka_three-stage_2021}. 
 Simultaneous with the high-precision clock interrogation, the clock laser will be stabilised via high-bandwidth interrogation of the atoms, 
 granting less stringent requirements for laser pre-stability. This removes the need for a large stabilisation cavity, 
 opening the door to portability and reducing the susceptibility of the system to mechanical shock. 
Important remarks:

\vspace{0.3cm}

% Uncomment the following if you wish to include a figure in the abstract
%
\begin{figure}[h]
\begin{center}
\includegraphics*[height=2.0cm]{schematic.png}
\end{center}
\caption{The caption for the figure should go here.}
\end{figure}

\setlength{\parskip}{\smallskipamount}
\noindent 

% Your references 

\printbibliography
% End your references

\end{document}
